% GENERATED FILE. Edit canonical lessons, then run npm run generate:books.
% canonical-source-hash: fnv1a64:0e8c90c1fbdf9193
% canonical-lessons: AR-C235-mathaf, AR-C235-malab, AR-C235-rasif, AR-C235-taqatu, AR-C235-sifara

\chapter{Museum, Stadium, Pavement, Crossroads, Embassy}
\label{ch:ar-museum-stadium-pavement-crossroads-embassy}
\hlchaptermodality{235}

\begin{chapteropening}
\textbf{By the end of this chapter:} \emph{I can say a museum, a stadium, a pavement, a crossroads and an embassy.}

Complete the last lesson of chapter 235: I can say a museum, a stadium, a pavement, a crossroads and an embassy.
\end{chapteropening}

\section[\texorpdfstring{\ar{متحف} (matḥaf)}{متحف (matḥaf)}]{\ar{متحف} (matḥaf) — a museum}

\label{lesson:AR-C235-mathaf}



\begin{quote}
Before the new one: say the Arabic for an alarm clock, then the Arabic for a calendar.
\end{quote}

\subsection*{You'll want to know: \ar{متحف}}
\textbf{\ar{متحف}} — \emph{matḥaf} — \textquotedblleft{}a museum\textquotedblright{}.

\begin{cousinweb}
\textbf{\ar{متحف}} (\emph{matḥaf}) is built on the native Arabic root \emph{t-ḥ-f}, \textquotedblleft{}a rare gift\textquotedblright{} (Lane). \emph{Matḥaf} is a modern place noun built on \emph{tuḥfa}, \textquotedblleft{}a rarity\textquotedblright{}: \textquotedblleft{}a place of rarities\textquotedblright{}.
\end{cousinweb}

\begin{grammarlens}[title={What you've built}]
One more word for this chapter.
\end{grammarlens}

\subsection*{Your turn}
\begin{itemize}
\raggedright
  \item \emph{Say it:} \emph{matḥaf}
  \item \emph{Say it:} \emph{matḥaf}, once more
  \item \emph{Say it:} say \emph{matḥaf}
  \item \emph{Recall:} say the Arabic for an alarm clock, then the Arabic for a calendar, then say \emph{matḥaf} again
\end{itemize}

\begin{tcolorbox}[breakable,colback=teal!4,colframe=teal!35!black,title={Before you move on}]
What does \ar{متحف} mean? (\textquotedblleft{}A museum\textquotedblright{}.) Say it once more.
\end{tcolorbox}
\section[\texorpdfstring{\ar{ملعب} (malʿab)}{ملعب (malʿab)}]{\ar{ملعب} (malʿab) — a stadium, a sports ground}

\label{lesson:AR-C235-malab}



\begin{quote}
Before the new one: say the Arabic for a calendar, then the Arabic for a museum.
\end{quote}

\subsection*{You'll want to know: \ar{ملعب}}
\textbf{\ar{ملعب}} — \emph{malʿab} — \textquotedblleft{}a stadium, a sports ground\textquotedblright{}.

\begin{cousinweb}
\textbf{\ar{ملعب}} (\emph{malʿab}) is built on the native Arabic root \emph{l-ʿ-b}, \textquotedblleft{}to play\textquotedblright{} (Lane). \emph{Malʿab} is the \emph{ma-} place noun, \textquotedblleft{}a place of play\textquotedblright{}.
\end{cousinweb}

\begin{grammarlens}[title={What you've built}]
One more word for this chapter.
\end{grammarlens}

\subsection*{Your turn}
\begin{itemize}
\raggedright
  \item \emph{Say it:} \emph{malʿab}
  \item \emph{Say it:} \emph{malʿab}, once more
  \item \emph{Say it:} say \emph{malʿab}
  \item \emph{Recall:} say the Arabic for a calendar, then the Arabic for a museum, then say \emph{malʿab} again
\end{itemize}

\begin{tcolorbox}[breakable,colback=teal!4,colframe=teal!35!black,title={Before you move on}]
What does \ar{ملعب} mean? (\textquotedblleft{}A stadium, a sports ground\textquotedblright{}.) Say it once more.
\end{tcolorbox}
\section[\texorpdfstring{\ar{رصيف} (raṣīf)}{رصيف (raṣīf)}]{\ar{رصيف} (raṣīf) — a pavement; a platform}

\label{lesson:AR-C235-rasif}



\begin{quote}
Before the new one: say the Arabic for a museum, then the Arabic for a stadium.
\end{quote}

\subsection*{You'll want to know: \ar{رصيف}}
\textbf{\ar{رصيف}} — \emph{raṣīf} — \textquotedblleft{}a pavement; a platform\textquotedblright{}.

\begin{cousinweb}
\textbf{\ar{رصيف}} (\emph{raṣīf}) is built on the native Arabic root \emph{r-ṣ-f}, \textquotedblleft{}to set stones close\textquotedblright{} (Lane). Lane's \emph{raṣīf} is stones set close together.
\end{cousinweb}

\begin{grammarlens}[title={What you've built}]
One more word for this chapter.
\end{grammarlens}

\subsection*{Your turn}
\begin{itemize}
\raggedright
  \item \emph{Say it:} \emph{raṣīf}
  \item \emph{Say it:} \emph{raṣīf}, once more
  \item \emph{Say it:} say \emph{raṣīf}
  \item \emph{Recall:} say the Arabic for a museum, then the Arabic for a stadium, then say \emph{raṣīf} again
\end{itemize}

\begin{tcolorbox}[breakable,colback=teal!4,colframe=teal!35!black,title={Before you move on}]
What does \ar{رصيف} mean? (\textquotedblleft{}A pavement; a platform\textquotedblright{}.) Say it once more.
\end{tcolorbox}
\section[\texorpdfstring{\ar{تقاطع} (taqāṭuʿ)}{تقاطع (taqāṭuʿ)}]{\ar{تقاطع} (taqāṭuʿ) — a crossroads, a junction}

\label{lesson:AR-C235-taqatu}



\begin{quote}
Before the new one: say the Arabic for a stadium, then the Arabic for a pavement; a platform.
\end{quote}

\subsection*{You'll want to know: \ar{تقاطع}}
\textbf{\ar{تقاطع}} — \emph{taqāṭuʿ} — \textquotedblleft{}a crossroads, a junction\textquotedblright{}.

\begin{cousinweb}
\textbf{\ar{تقاطع}} (\emph{taqāṭuʿ}) is built on the native Arabic root \emph{q-ṭ-ʿ}, \textquotedblleft{}to cut, to sever\textquotedblright{} (Lane). \emph{Taqāṭuʿ} is the verbal noun of Form VI, \textquotedblleft{}cutting across each other\textquotedblright{}.
\end{cousinweb}

\begin{grammarlens}[title={What you've built}]
One more word for this chapter.
\end{grammarlens}

\subsection*{Your turn}
\begin{itemize}
\raggedright
  \item \emph{Say it:} \emph{taqāṭuʿ}
  \item \emph{Say it:} \emph{taqāṭuʿ}, once more
  \item \emph{Say it:} say \emph{taqāṭuʿ}
  \item \emph{Recall:} say the Arabic for a stadium, then the Arabic for a pavement; a platform, then say \emph{taqāṭuʿ} again
\end{itemize}

\begin{tcolorbox}[breakable,colback=teal!4,colframe=teal!35!black,title={Before you move on}]
What does \ar{تقاطع} mean? (\textquotedblleft{}A crossroads, a junction\textquotedblright{}.) Say it once more.
\end{tcolorbox}
\section[\texorpdfstring{\ar{سفارة} (sifāra)}{سفارة (sifāra)}]{\ar{سفارة} (sifāra) — an embassy}

\label{lesson:AR-C235-sifara}



\begin{quote}
Before the new one: say the Arabic for a pavement; a platform, then the Arabic for a crossroads.
\end{quote}

\subsection*{You'll want to know: \ar{سفارة}}
\textbf{\ar{سفارة}} — \emph{sifāra} — \textquotedblleft{}an embassy\textquotedblright{}.

\begin{cousinweb}
\textbf{\ar{سفارة}} (\emph{sifāra}) is built on the native Arabic root \emph{s-f-r}, \textquotedblleft{}to sweep, to uncover\textquotedblright{} (Lane), which also gives \emph{safara bayna}, \textquotedblleft{}to make peace between\textquotedblright{}. \emph{Sifāra} is the office of \emph{safīr}, \textquotedblleft{}a go-between\textquotedblright{}; the embassy is a modern sense.
\end{cousinweb}

\begin{grammarlens}[title={What you've built}]
One more word for this chapter.
\end{grammarlens}

\subsection*{Your turn}
\begin{itemize}
\raggedright
  \item \emph{Say it:} \emph{sifāra}
  \item \emph{Say it:} \emph{sifāra}, once more
  \item \emph{Say it:} say \emph{sifāra}
  \item \emph{Recall:} say the Arabic for a pavement; a platform, then the Arabic for a crossroads, then say \emph{sifāra} again
\end{itemize}

\begin{tcolorbox}[breakable,colback=teal!4,colframe=teal!35!black,title={Before you move on}]
What does \ar{سفارة} mean? (\textquotedblleft{}An embassy\textquotedblright{}.) Say it once more.
\end{tcolorbox}
